\clearpage
\section{Introduction}\label{section::introduction}

\subsection{Research Background}
Over the past decade, the field of computer vision has witnessed substantial advances in core tasks such as object detection, image segmentation, and multi-object tracking. The rise of deep learning has been the primary driver of this progress: its powerful nonlinear modeling capacity enables algorithms to autonomously learn complex visual patterns from large-scale data, thereby achieving accurate recognition and localization of objects in images and videos \cite{113,28_kaur_comprehensive_2023}.

Among a broad range of computer vision applications, multi-object tracking (MOT) stands out as a particularly challenging and consequential research topic \cite{104}. The essence of MOT goes beyond accurately detecting all objects in a scene; it further requires assigning and maintaining a unique identity for each object across consecutive video frames \cite{103}. MOT is a core component of critical domains including autonomous driving \cite{117}, intelligent surveillance \cite{59_wang_deep_2024}, robotics \cite{188}, intelligent traffic management \cite{144}, and crowd analytics \cite{158}. In these dynamic and complex real-world environments, real-time and accurate detection and tracking are foundational to ensuring system reliability and effectiveness.

At present, the tracking-by-detection paradigm has become the de facto solution in industrial MOT systems due to its robustness and flexibility, and its overall performance hinges on the accuracy and efficiency of the underlying object detector \cite{132}. Among modern detectors, the You Only Look Once (YOLO) family has been widely adopted for real-time applications owing to its excellent balance between speed and accuracy \cite{02_noauthor_review_nodate,141}. YOLO formulates object detection as a single-stage regression problem that directly predicts bounding boxes and class probabilities, enabling high-throughput processing. Since its inception, the YOLO series has undergone multiple iterations, and this continual evolution has laid a solid foundation for improving the end-to-end performance of MOT systems.

On the tracking side, DeepSORT is a significant extension of the Simple Online and Realtime Tracking (SORT) algorithm that integrates deep appearance features through a learned embedding model. By fusing motion dynamics with visual appearance, DeepSORT is able to better handle long-term occlusions and re-identification, substantially reducing identity switches (IDSW) \cite{211}. Thanks to its high tracking accuracy, computational efficiency, and seamless integration with various detectors, DeepSORT has demonstrated strong practical value in domains such as traffic surveillance \cite{138}, warehousing and logistics \cite{212}, sports analytics \cite{213}, and crowd monitoring \cite{120}.

\subsection{Challenges}
Despite the significant advances driven by deep learning, real-time MOT in complex, high-density environments—particularly for crowd analytics—still faces a number of severe challenges rooted in the inherent complexity of real-world scenes. These include drastic fluctuations in crowd density, severe inter-person occlusions, highly variable illumination, and appearance changes induced by viewpoint and scale variations \cite{30_khan_passenger_2020,01_noauthor_deep_2024,66_zhang_bus_2020}.

First, high density and occlusion are pervasive challenges in crowd analysis. In crowded scenes, frequent mutual occlusions lead to blurred object boundaries, making precise separation and counting difficult. The severity of occlusions typically increases with crowd density, substantially complicating counting and tracking\cite{33or34_li_lightweight_2023}\cite{30_khan_passenger_2020}\cite{01_noauthor_deep_2024}.

Second, large variations in scale and viewpoint pose another major challenge. Due to perspective effects, individuals exhibit markedly different pixel scales within the same image—small in the far field and large in the near field. Such scale disparities hinder methods that rely on fixed-size assumptions for detection or segmentation. In addition, differences in camera viewpoints alter spatial configuration and apparent geometry, further exacerbating the difficulty of accurate counting\cite{29_kaur_systematic_2024}\cite{30_khan_passenger_2020}.

Third, complex backgrounds and varying illumination introduce significant interference. Real-world scenes often contain textured background elements such as buildings, vehicles, and vegetation, whose appearance may be confused with human targets, leading to false positives. Meanwhile, drastic illumination changes—strong shadows, over-exposure, or low-light conditions—can degrade image quality and reduce target discriminability, impeding accurate recognition\cite{33or34_li_lightweight_2023}\cite{30_khan_passenger_2020}.

Finally, substantial intra-class appearance variation further increases complexity. Large differences in clothing, pose, and motion across individuals\cite{33or34_li_lightweight_2023}\cite{66_zhang_bus_2020}\cite{30_khan_passenger_2020}\cite{01_noauthor_deep_2024} place stringent demands on the generalization capacity of counting and tracking models. These factors often co-occur and interact, making accurate crowd counting a problem that requires advanced and robust algorithms.

\subsection{Motivation and Objectives}
Motivated by the above challenges, we focus on railway platforms—public transport environments characterized by highly dense crowds and frequent occlusions—where accurate, real-time detection, tracking, and counting are of pressing practical importance and research value. Reliable crowd management is not only critical to ensuring operational safety\cite{214}, but also fundamental for optimizing passenger flow\cite{215} and improving emergency response efficiency\cite{216}.

Accordingly, our goal is to develop an efficient and robust automated crowd counting system tailored to the railway platform scenario. Leveraging state-of-the-art deep learning techniques, the system aims to effectively address high-density and severe occlusion, providing reliable data support for platform management, operational efficiency, and passenger experience.

To this end, we adopt YOLOv11—an important iteration in the YOLO family—as the object detector, and integrate DeepSORT as the tracker. YOLOv11 is selected for its strong balance between speed and accuracy\cite{141}, and for architectural advances that are expected to improve small-object detection and robustness in complex scenes. DeepSORT is chosen for its high tracking accuracy, computational efficiency, and compatibility with diverse detectors, which make it well-suited to handling occlusion and re-identification. By tightly coupling YOLOv11's strong detection capacity with DeepSORT's robust association mechanism, we seek to build a system that delivers high accuracy, real-time performance, and robustness for detection, tracking, and counting in dense and dynamic railway platform environments—ultimately supporting intelligent transportation and public safety.

\subsection{Main Contributions}
We focus on the specific characteristics of railway platforms—high crowd density, frequent occlusions under fixed viewpoints, and stringent requirements on both real-time performance and accuracy—and explore a YOLOv11-based approach to crowd detection, tracking, and counting. Through comprehensive evaluation on publicly available datasets and a self-constructed platform dataset, we assess the applicability of YOLOv11 to this scenario and provide empirical evidence to inform future research and deployment.

Our main contributions are summarized as follows:\\
First, we apply the state-of-the-art YOLOv11 architecture to dense crowd counting. While the YOLO family has achieved notable success in generic object detection, directly deploying its important iteration version for crowd counting and conducting an in-depth performance analysis remains a novel and informative endeavor.\\

Second, we propose two improved variants of DeepSORT and train a lightweight convolutional neural network (CNN), EfficientNet-B0, for specialized head feature extraction to enhance robustness under occlusions.\\

Third, to train and rigorously evaluate the proposed models, we employ multiple public datasets—CrowdHuman\cite{shao2018crowdhuman}, NWPU-Crowd\cite{nwpcrowd}, Open Sensor Data for Rail 2023\cite{OpenSensorDataforRail2023}, and RailEye3D\cite{RailEye3D_Dataset}—together with a self-constructed railway platform dataset named MOT RailwayPlatformCrowdHead (MOT-RPCH). This comprehensive dataset contains 27 video sequences with 24,788 frames, 89,087 bounding boxes, and 885 unique human head identities. For MOT evaluation, we selected a representative subset of 20 video sequences with 18,548 frames, 73,799 bounding boxes, and 647 identities (\textbf{Tab.}~\ref{tab:20MOTValVideos}), providing a valuable benchmark for railway platform crowd analytics.\\

Finally, we conduct a comprehensive comparison against a standard 8-dimensional Kalman filtering baseline. Experimental results clearly demonstrate the superiority of our approach and validate the feasibility and effectiveness of combining YOLOv11 with improved DeepSORT for crowd analytics.

\subsection{Outline of the Dissertation}
The remainder of this paper is organized as follows:
\begin{itemize}
  \item \textbf{Chapter 2: Related Work} reviews mainstream methods, key techniques, and representative systems in object detection and MOT, and clarifies the relationship and distinctions between our work and prior art.
  \item \textbf{Chapter 3: Data} describes the public and self-constructed datasets used in this study, including their sources and salient characteristics, and discusses their suitability for crowd analytics on railway platforms.
  \item \textbf{Chapter 4: Methods} details the core models used in our study, including the principles, architectures, and key components of YOLOv11 (detector) and DeepSORT (tracker). The chapter presents our detect–track–count pipeline with three Kalman state-space configurations (CV-8D, CA-12D, and Phys-6D), introduces the scene and imaging model with 2D planar trajectory visualization, and defines comprehensive MOT and counting metrics for evaluation.
  \item \textbf{Chapter 5: Experiments and Results} presents the experimental setup, implementation details, and comprehensive quantitative and qualitative analyses. The evaluation metrics defined in \textbf{Chapter} 4 are used to quantify performance.
  \item \textbf{Chapter 6: Discussion} analyzes the strengths, limitations, and potential implications of the proposed approach for detection, tracking, and counting on railway platforms.
  \item \textbf{Chapter 7: Conclusion and Future Work} summarizes the main findings and contributions, and outlines directions for future research.
\end{itemize}

